% !TeX program = pdflatex
% Kompilieren mit: pdflatex rfn_aufbau.tex (zweimal)
% Gleicher Stil wie slides.tex (Koppeltest). Bilder liegen in img/.
% img/JLU_Giessen-Logo.pdf aus dem Ordner der anderen Praesentation hierher kopieren.
\documentclass[10pt,aspectratio=169]{beamer}

% ── Theme ──────────────────────────────────────────────────────────────────
\usetheme{metropolis}

% ── Pakete ─────────────────────────────────────────────────────────────────
\usepackage[utf8]{inputenc}
\usepackage[T1]{fontenc}
\usepackage{lmodern}
\usepackage[ngerman]{babel}
\usepackage{tikz}
\usetikzlibrary{positioning,arrows.meta,fit,backgrounds,calc,shapes.geometric}
\usepackage{booktabs}
\usepackage{array}
\newcolumntype{L}[1]{>{\raggedright\arraybackslash}p{#1}}
\usepackage{siunitx}
\sisetup{locale = DE, per-mode=symbol}
\usepackage{textcomp}
\makeatletter
\let\jlu@olditemize\itemize
\let\jlu@oldenditemize\enditemize
\renewenvironment{itemize}
  {\jlu@olditemize\setlength{\itemsep}{1pt}\setlength{\parsep}{0pt}\setlength{\topsep}{2pt}}
  {\jlu@oldenditemize}
\makeatother

% ── Farben ─────────────────────────────────────────────────────────────────
\definecolor{giessenblau}{RGB}{28,54,88}
\definecolor{hellblau}{RGB}{230,234,240}
\definecolor{grauton}{RGB}{95,102,110}
\definecolor{hellgrau}{RGB}{236,236,237}

\setbeamercolor{palette primary}{bg=giessenblau,fg=white}
\setbeamercolor{frametitle}{bg=white,fg=giessenblau}
\setbeamercolor{progress bar}{fg=giessenblau,bg=giessenblau!20}
\setbeamercolor{alerted text}{fg=giessenblau}
\setbeamercolor{example text}{fg=grauton}
\setbeamertemplate{itemize item}{\color{giessenblau}\textbullet}
\setbeamertemplate{itemize subitem}{\color{grauton}\textendash}
\setbeamercolor{block title}{bg=giessenblau,fg=white}
\setbeamercolor{block body}{bg=hellblau}
\setbeamercolor{block title alerted}{bg=grauton,fg=white}
\setbeamercolor{block body alerted}{bg=hellgrau,fg=black}

\metroset{sectionpage=none, numbering=fraction, progressbar=frametitle}
\makeatletter
\setlength{\metropolis@progressinheadfoot@linewidth}{2pt}
\makeatother

\newcommand{\todo}[1]{\textcolor{giessenblau}{\bfseries [TODO: #1]}}
\newcommand{\bildplatzhalter}[3]{%
  \begin{tikzpicture}
    \draw[grauton, dashed, thick, rounded corners=2pt] (0,0) rectangle (#1,#2);
    \node[grauton, align=center, text width=#1 cm-0.8cm] at ({#1/2},{#2/2}) {\small #3};
  \end{tikzpicture}%
}
\newcommand{\hinweis}[1]{{\small\color{grauton} #1}}
% Pro/Contra-Zeichen
\newcommand{\pro}{\textcolor{giessenblau}{\textbf{+}}}
\newcommand{\con}{\textcolor{grauton}{\textbf{--}}}

% ── Meta ───────────────────────────────────────────────────────────────────
\newcommand{\vortragsdatum}{22.10.2026}
\title{Neue Blendenaufnahme des HF-Neutralisators}
\subtitle{Aufbau, Dichtkonzept und Bauteile}
\author{Konrad Wehkamp}
\date{}

\begin{document}

% ============================================================
{
\setbeamertemplate{footline}{}
\begin{frame}[plain,c]
  \begin{tikzpicture}[remember picture, overlay]
    \node[anchor=north west, inner sep=0pt] at ([xshift=0.9cm,yshift=-0.6cm]current page.north west)
      {\includegraphics[height=1.0cm]{img/JLU_Giessen-Logo.pdf}};
    \node[anchor=north east, inner sep=0pt, text width=6cm, align=right, font=\scriptsize\color{grauton}]
      at ([xshift=-0.9cm,yshift=-0.55cm]current page.north east)
      {I.\ Physikalisches Institut\\AG Ionentriebwerke};
    \draw[giessenblau, line width=0.8pt] ([yshift=-1.7cm]current page.north west) -- ([yshift=-1.7cm]current page.north east);
    \draw[giessenblau, line width=0.8pt] ([yshift=1.6cm]current page.south west) -- ([yshift=1.6cm]current page.south east);
  \end{tikzpicture}

  \begin{center}
    {\usebeamerfont{title}\LARGE\bfseries\color{giessenblau} Neue Blendenaufnahme des HF-Neutralisators}\\[0.7em]
    {\Large\color{grauton} Aufbau, Dichtkonzept und Bauteile}\par
    \vspace{0.9cm}
    {\large\textbf{Konrad Wehkamp}}\\[0.3em]
    {\small Masterstudiengang PTRA (Physics and Technology for Space Applications)}\par
    \vspace{0.6cm}
    {\small
      \begin{tabular}{r l}
        Betreuung: & Jana Zorn\\
        Rahmen: & Spezialisierungsprojekt / Vorbereitung Masterarbeit \\
        Datum: & \vortragsdatum \\
      \end{tabular}
    }
  \end{center}
\end{frame}
}

% ============================================================
\begin{frame}{Warum eine neue Aufnahme}
  \vfill
  Der erste Neutralisator läuft. Die untere Aufnahme hat aber drei Schwächen,
  die bei der Charakterisierung stören.

  \vspace{0.6em}
  \begin{itemize}
    \item Die Blende liegt lose auf einem GFK-Ring. Ob ihr Rand am Glas anliegt, entscheidet die Toleranz,
          ein Teil des Gases kann an den Löchern vorbei in den Tank
    \item Der O-Ring dichtet axial unter dem Glasboden, seine Verpressung hängt an der Höhe des Glasfußes
    \item Kontaktiert wird die Blende von der Austrittsseite über ein Gewinde
  \end{itemize}

  \vspace{0.6em}
  Gleichzeitig soll die Blende ein Wechselteil werden, um verschiedene Lochbilder zu testen.
  \vfill
\end{frame}

% ============================================================
\begin{frame}{Anforderungen}
  \vfill
  \begin{columns}[T]
    \begin{column}{0.5\textwidth}
      \small
      \begin{block}{Funktion}
        \begin{itemize}
          \item Glasraum gasdicht gegen die Spulenkammer
          \item Blende fest angedrückt, kein Gas am Lochfeld vorbei
          \item Kein zusätzliches Material hinter der Blende, auch nicht weiter außen
          \item Blende isoliert und von oben kontaktiert
        \end{itemize}
      \end{block}
    \end{column}
    \begin{column}{0.46\textwidth}
      \small
      \begin{block}{Handhabung}
        \begin{itemize}
          \item Blende tauschbar, Lochbilder bis Ø\,40\,mm
          \item RFN komplett montiert an die Frontplatte schrauben und wieder abnehmen
          \item Glasfuß bleibt die einzige feste Einspannung des Glases
          \item Einziger harter Anschlag: oberer GFK-Ring auf dem Edelstahlring
          \item Vakuumtauglich, keine eingeschlossenen Volumen
        \end{itemize}
      \end{block}
    \end{column}
  \end{columns}
  \vfill
\end{frame}

% ============================================================
\begin{frame}{Gesamtbild}
  \centering
  \includegraphics[height=5.5cm]{img/rfn_schnitt3d.png}

  \vspace{0.4em}
  \hinweis{Viertelschnitt durch Frontplatte, Träger, Edelstahlring, GFK-Ring oben und Glas.
  Rechts ist eine Federschraube angeschnitten, dahinter der Kontaktbolzen.}
\end{frame}

% ============================================================
\begin{frame}{Halbschnitt}
  \centering
  \includegraphics[width=\textwidth,height=5.8cm,keepaspectratio]{img/halbschnitt.pdf}

  \hinweis{Maße in mm, z ab Blendenunterseite. Hinter der Blende nur noch 1,5\,mm Klemmring statt 1,9\,mm.}
\end{frame}

% ============================================================
\begin{frame}{Bauteile}
  \begin{columns}[T]
    \begin{column}{0.66\textwidth}
      \scriptsize
      \renewcommand{\arraystretch}{1.15}
      \begin{tabular}{@{}L{2.1cm}L{1.6cm}L{4.6cm}@{}}
        \toprule
        \textbf{Teil} & \textbf{Werkstoff} & \textbf{Aufgabe} \\
        \midrule
        Blende & Mo & 0,5\,mm, Wechselteil, Ø\,72, Lochfeld bis Ø\,40 \\
        Klemmring & 1.4301 & 1,5\,mm, klemmt die Blende, Lasche für den Kontakt \\
        Träger & GFK & Ø\,80 g6 in der Frontplatte, isoliert, Gewindeeinsätze \\
        Edelstahlring & 1.4301 & beide O-Ring-Nuten, Gewinde für Federn und Bleche \\
        GFK-Ring oben & GFK & drückt federnd auf den Glasfuß \\
        O-Ring radial & FFKM & 61,5\,×\,2,65 nach DIN\,3771, auf dem Glasfuß \\
        O-Ring flach & FFKM & 64\,×\,1, Fuge Edelstahlring / Träger \\
        Keramikhülse & Al\textsubscript{2}O\textsubscript{3} & führt den Kontaktbolzen isoliert nach oben \\
        Federn, Hülsen & 1.4310, A2 & 9 Druckfedern auf Distanzhülsen, Scheibe darunter \\
        \bottomrule
      \end{tabular}
    \end{column}
    \begin{column}{0.31\textwidth}
      \centering
      \includegraphics[height=6.2cm]{img/rfn_explosion.png}
    \end{column}
  \end{columns}
\end{frame}

% ============================================================
\begin{frame}{Kraftfluss}
  \vfill
  \centering
  \begin{tikzpicture}[
      box/.style={draw=giessenblau, fill=hellblau, rounded corners=2pt, minimum height=0.8cm,
                  minimum width=2.1cm, font=\footnotesize, align=center},
      spr/.style={box, draw=grauton, fill=hellgrau},
      ar/.style={-{Latex[length=2mm]}, giessenblau, thick}]
    \node[box] (go) {GFK-Ring\\oben};
    \node[box, right=0.5cm of go] (gl) {Glasfuß};
    \node[box, right=0.5cm of gl] (bl) {Blende};
    \node[box, right=0.5cm of bl] (kr) {Klemmring};
    \node[box, below=1.0cm of kr] (tr) {Träger\\6\,× M2};
    \node[box, left=0.5cm of tr] (st) {Edelstahlring\\3\,× M2,5};
    \node[box, left=0.5cm of st] (dh) {9\,× M3,\\Distanzhülse};
    \node[spr, left=0.5cm of dh] (fe) {9\,× Feder};
    \draw[ar] (go) -- (gl); \draw[ar] (gl) -- (bl); \draw[ar] (bl) -- (kr);
    \draw[ar] (kr) -- (tr); \draw[ar] (tr) -- (st); \draw[ar] (st) -- (dh);
    \draw[ar] (dh) -- (fe); \draw[ar] (fe) -- (go);
  \end{tikzpicture}

  \vspace{1.0em}
  \raggedright
  \begin{itemize}
    \item Im geschlossenen Kreis sitzt genau ein nachgiebiges Element, die neun Federn. Rund 110\,N halten das Glas auf der Blende
    \item Die Distanzhülsen begrenzen den Federweg. Die Kraft hängt nicht davon ab, wie fest man die Schrauben anzieht
    \item Zwischen GFK-Ring und Stahl bleiben 0,3\,mm Luft. Erst bei Überlast liegt er dort an
    \item Die Frontplatte liegt außerhalb des Kreises. Sie trägt den RFN, klemmt aber nicht das Glas
  \end{itemize}
  \vfill
\end{frame}

% ============================================================
\begin{frame}{Dichtkonzept}
  \begin{columns}[T]
    \begin{column}{0.47\textwidth}
      \small
      \textbf{Glasraum gegen Spulenkammer}
      \begin{itemize}
        \item Radialer O-Ring auf dem Glasfuß. Die Höhe des Glases spielt keine Rolle mehr
        \item Verpressung 19\,\%, bei ganz außermittigem Glas 15 bis 23\,\%
        \item Sitzt im Edelstahlring, 5\,mm über der Blende
        \item Zweiter, flacher O-Ring schließt die Fuge zwischen Stahl und Träger
      \end{itemize}
      \vspace{0.4em}
      \textbf{Glasraum gegen Tank, am Lochfeld vorbei}
      \begin{itemize}
        \item Glasboden auf Blende und Blende auf Träger, zwei Anpressflächen hintereinander
        \item Nicht elastisch dicht, aber nicht mehr lose
      \end{itemize}
    \end{column}
    \begin{column}{0.5\textwidth}
      \centering
      \includegraphics[width=\linewidth]{img/rfn_detail.png}
    \end{column}
  \end{columns}
\end{frame}

% ============================================================
\begin{frame}{Kontaktierung der Blende}
  \begin{columns}[T]
    \begin{column}{0.45\textwidth}
      \small
      \vspace{0.6em}
      Eine Gewindestange M2,5 sitzt in der Lasche des Klemmrings und läuft in einer
      Keramikhülse durch Edelstahlring und GFK-Ring nach oben.

      \vspace{0.6em}
      Die Mutter oben stützt sich auf die Hülse. Der Bolzen zieht so die Lasche an den Träger,
      ohne den federnden GFK-Ring zu belasten.

      \vspace{0.6em}
      Zum Blendenwechsel Kontermutter, Kabelschuh und Mutter abnehmen und den Klemmring
      samt Bolzen nach hinten herausziehen. Glas und GFK-Ring bleiben montiert.
    \end{column}
    \begin{column}{0.52\textwidth}
      \centering
      \includegraphics[width=\linewidth,height=6.0cm,keepaspectratio]{img/kontakt.pdf}
    \end{column}
  \end{columns}
\end{frame}

% ============================================================
\begin{frame}{Montage und Befestigung}
  \begin{columns}[T]
    \begin{column}{0.5\textwidth}
      \footnotesize
      \begin{enumerate}
        \item Träger mit Flach-O-Ring an den Edelstahlring schrauben
        \item Blende und Klemmring von hinten einsetzen, Kontaktbolzen durch die Hülse
        \item Glas durch den radialen O-Ring bis auf die Blende schieben
        \item GFK-Ring auflegen, Hülsen, Federn und neun M3 einsetzen
        \item Fertigen RFN in die Frontplatte stecken, auf den Passstift drehen
        \item Drei M2 durch offene Nuten im GFK-Ring anziehen
      \end{enumerate}
      \vspace{0.4em}
      \hinweis{Abnehmen geht umgekehrt, der RFN bleibt dabei komplett montiert.}
    \end{column}
    \begin{column}{0.48\textwidth}
      \centering
      \includegraphics[width=\linewidth,height=6.0cm,keepaspectratio]{img/draufsicht.pdf}
    \end{column}
  \end{columns}
\end{frame}

% ============================================================
\begin{frame}{Vakuum: eingeschlossene Volumen}
  \begin{columns}[T]
    \begin{column}{0.44\textwidth}
      \footnotesize
      Rundum verpresste Sacklöcher und Fugen wären virtuelle Lecks.
      \begin{itemize}
        \item Ringnut unten im Stahlring verbindet M3-Gewinde, M2,5-Löcher und Hülsenbohrung
        \item Drei Bohrungen Ø\,1 im Träger führen nach hinten in den Tank
        \item Kerben an Flachnut, Klemmring und unter der Keramikhülse
      \end{itemize}
      \vspace{0.4em}
      Der Ringraum um den Glasfuß (0,6\,cm\textsuperscript{3}) bleibt absichtlich zu, er liegt zwischen den Anpressflächen an der Blende. Gegenüber dem Gasfluss ist er bedeutungslos.
    \end{column}
    \begin{column}{0.54\textwidth}
      \centering
      \includegraphics[width=\linewidth,height=5.6cm,keepaspectratio]{img/entlueftung.pdf}
    \end{column}
  \end{columns}
\end{frame}

% ============================================================
\begin{frame}{Vor- und Nachteile}
  \begin{columns}[T]
    \begin{column}{0.49\textwidth}
      \footnotesize
      \begin{block}{Vorteile}
        \begin{itemize}
          \item[\pro] Blende definiert angedrückt, Nebenweg zweifach gepresst statt lose
          \item[\pro] O-Ring unabhängig von der Glashöhe
          \item[\pro] Federkraft definiert, Toleranz des Glasfußes ändert sie nur um rund 1\,N je Feder
          \item[\pro] 1,5 statt 1,9\,mm Material hinter der Blende
          \item[\pro] Blendenwechsel ohne Glas oder GFK-Ring abzunehmen
          \item[\pro] RFN montiert an- und abschraubbar
          \item[\pro] Komplett parametrisch, 160 Parameter mit Prüfungen
        \end{itemize}
      \end{block}
    \end{column}
    \begin{column}{0.49\textwidth}
      \footnotesize
      \begin{alertblock}{Nachteile und Risiken}
        \begin{itemize}
          \item[\con] Mehr Einzelteile und eine zweite Dichtstelle
          \item[\con] FFKM-O-Ring am Glasfuß wie im ersten RFN, Temperatur am neuen Ort noch nicht gerechnet
          \item[\con] Träger aus GFK liegt an der Blende und trägt neun Gewindeeinsätze
          \item[\con] Nur 0,1\,mm Spiel am Glasfuß, Durchmesser und Rundheit müssen stimmen
          \item[\con] Dichtheit am Lochfeld hängt an der Planheit von Glasboden und Blende
          \item[\con] Abdeckbleche brauchen neue Bohrungen
        \end{itemize}
      \end{alertblock}
    \end{column}
  \end{columns}
\end{frame}

% ============================================================
\begin{frame}{Offene Punkte}
  \vfill
  \begin{itemize}
    \item Glasfuß vermessen: Durchmesser, Rundheit, Höhe. Nuten und Bohrungen folgen automatisch
    \item COMSOL mit der neuen Geometrie: Temperatur an Träger, O-Ring und Federn.
          Reserve bei zu heißem Träger: Polyimid (Vespel)
    \item Druckfedern aus dem Katalog wählen, die Werte sind noch Platzhalter. Prüfen, ob über dem GFK-Ring 2\,mm Platz für die Schraubenköpfe sind
    \item Abdeckbleche an die versetzten Gewinde anpassen
    \item Frontplatte und Teile in die Fertigung geben
  \end{itemize}
  \vfill
\end{frame}

% ============================================================
{
\setbeamertemplate{footline}{}
\begin{frame}[plain,c]
  \begin{tikzpicture}[remember picture, overlay]
    \node[anchor=north west, inner sep=0pt] at ([xshift=0.9cm,yshift=-0.6cm]current page.north west)
      {\includegraphics[height=1.0cm]{img/JLU_Giessen-Logo.pdf}};
    \draw[giessenblau, line width=0.8pt] ([yshift=-1.7cm]current page.north west) -- ([yshift=-1.7cm]current page.north east);
    \draw[giessenblau, line width=0.8pt] ([yshift=1.6cm]current page.south west) -- ([yshift=1.6cm]current page.south east);
  \end{tikzpicture}
  \begin{center}
    {\LARGE\bfseries\color{giessenblau} Vielen Dank}\\[0.8em]
    {\Large\color{grauton} Fragen?}\\[2.2em]
    {\small\color{grauton} Konrad Wehkamp \quad\textperiodcentered\quad Blendenaufnahme RFN \quad\textperiodcentered\quad \vortragsdatum}
  \end{center}
\end{frame}
}

\end{document}
